% =============================================================================
% main.tex -- Phase E2 manuscript SKELETON. NOT a draft.
% =============================================================================
% Governed by, in this order of authority: research/PAPER_CONTRACT.md (binding
% anti-drift contract) > research/CONTRIBUTION_LOCK.md (locked contribution,
% Formulation #2) > research/MANUSCRIPT_ARCHITECTURE.md (E1 architecture, which
% this file implements section-for-section).
%
% Every \subsection below carries a compact "draftnote" block
% (PURPOSE / QUESTION / ALLOWED EVIDENCE / KEY POINTS / MUST NOT SAY) plus an
% % EVIDENCE: comment anchor. E3 drafting must replace the draftnote content
% with real prose that satisfies exactly what the block specifies -- it must
% not invent claims, numbers, or citations beyond what the block licenses.
%
% TITLE DECISION (per the E2 brief -- record in main.tex and the final report):
% "Mechanism" was chosen over "classifier". Every locked evidence document
% (research/CONTRIBUTION_LOCK.md, research/contribution_lock.csv,
% research/PAPER_CONTRACT.md, research/MANUSCRIPT_ARCHITECTURE.md,
% research/EXPERIMENT_1_DATA_DICTIONARY.md) exclusively calls the two compared
% units "mechanisms" (rules_only / retrieval_only; "empirical_winner" is a
% mechanism, never a "classifier"). "Classifier" is not used anywhere in the
% frozen evidence for this comparison and is technically looser: the rules
% mechanism is an exact-match lookup, not a trained classifier, and calling it
% one would be a small but avoidable terminology drift the moment a reviewer
% checks the Experimental Design section against the word "classifier" in the
% title. No third conceptual framing was introduced.
%
% LATEX FORMAT (E0 decision, research/MANUSCRIPT_FORMAT_RESEARCH.md Sec.2.3):
% generic article class, pdfLaTeX-compatible, natbib + BibTeX, no venue-
% specific template. Packages below are the minimum needed to compile this
% skeleton (amsmath/amssymb for the three already-frozen equations E1-E3;
% natbib for \citep) -- nothing added "for later."
% =============================================================================

\documentclass[11pt]{article}

\usepackage{amsmath}
\usepackage{amssymb}
\usepackage[round]{natbib}

\newenvironment{draftnote}{%
  \par\noindent\begin{quote}\small\ttfamily
}{%
  \end{quote}\par
}

\title{Historical Consistency Predicts Mechanism Accuracy, Not Mechanism
Ranking: Evidence from a Controlled Synthetic Study}
% TODO: title placeholder pending final human sign-off among the five
% candidates ranked in research/MANUSCRIPT_ARCHITECTURE.md Sec.1. This is
% candidate #1 (the E1-recommended, most defensible option), with "classifier"
% swapped to "mechanism" per the decision recorded above. Do not substitute a
% third framing without updating research/MANUSCRIPT_ARCHITECTURE.md first.

\author{TODO: Author Name \\ \small TODO: Affiliation / independent researcher status \\ \small TODO: ORCID iD}
% AFFILIATION / AUTHOR PLACEHOLDER -- per the E2 brief, left as an explicit
% placeholder rather than pre-filled. research/MANUSCRIPT_FORMAT_RESEARCH.md
% Sec.1.12/1.14 has open items (ORCID reuse, endorsement) that must be resolved
% before this placeholder is finalized, not guessed here.

\date{TODO: submission date}

\begin{document}

\maketitle

\begin{abstract}
% =============================================================================
% ABSTRACT PLACEHOLDER -- DO NOT WRITE THE ABSTRACT IN E2. Structured TODO
% block per the E2 brief. Content requirements are exhaustive: the final
% abstract must contain exactly [A]-[H] below, in this order, and nothing else
% (research/MANUSCRIPT_ARCHITECTURE.md Sec.2).
% =============================================================================
\begin{draftnote}
TODO -- ABSTRACT, structured content requirements (draft at E3, not now):\\[2pt]
{[}A{]} PROBLEM: classification systems are usually built by choosing a
mechanism (rules, retrieval, a model) up front, without a measured signal for
whether that choice fits the data.\\
{[}B{]} RESEARCH QUESTION: can historical decision consistency, measured
before deployment, be used to select between qualitatively different
classification mechanisms? (verbatim scope: CONTRIBUTION\_LOCK.md \S11.A)\\
{[}C{]} METHOD: a pre-registered, 240-condition synthetic factorial experiment
(20 seeds $\times$ 6 nominal-consistency targets $\times$ 2 lexical
conditions), comparing exact-match rules against fuzzy retrieval.\\
{[}D{]} EXPERIMENTAL SETUP: one compressed sentence naming the two mechanisms
and the two lexical conditions (CLEAN / VARIED) -- not the full configuration
table (that belongs in Sec.~4, Table~T3).\\
{[}E{]} MAIN RESULT (positive half, 6a): realized ADS strongly correlated with
each mechanism's own accuracy (Pearson $r\approx0.909$--$0.959$ rules,
$0.948$--$0.955$ retrieval, both lexical conditions).\\
{[}F{]} LIMITING/NEGATIVE RESULT (6b): ADS did \emph{not} predict which
mechanism wins -- 100\% rule/empirical-winner agreement in the realized
0.70--0.90 ADS band vs.\ 0\% in the realized $\geq$0.90 band, because the
empirical winner is a constant function of the (separately manipulated)
lexical condition, unconditional on ADS across its full observed range.\\
{[}G{]} SCOPE LIMITATION: one synthetic generator, one lexical-perturbation
model, one motivating (non-evidentiary) production case study -- not a
deployment or generalization claim.\\
{[}H{]} IMPLICATION: a design-time selector built on consistency alone should
not be trusted to pick a mechanism without also accounting for representation
stability -- named as a direction, not something built or tested here.\\[4pt]
NUMBERS LIKELY IN THE FINAL ABSTRACT: the two Pearson-$r$ ranges (E) and the
100\%-vs-0\% band split (F), stated together, never separately. NUMBERS THAT
MUST STAY IN THE BODY ONLY: the flat 64.0\% aggregate agreement rate (32/50),
its Wilson CI [50.14\%, 75.86\%] and $p=0.0649$, the 120/120-vs-120/120
winner-constancy raw counts, and every production case-study number (91.2\%,
0.847, 0.964, 0.695) -- see research/MANUSCRIPT\_ARCHITECTURE.md \S2's table
for the full rationale on each exclusion.\\[4pt]
MUST NOT SAY: ``ADS selects the correct mechanism''; any claim implying
production data validates this experiment; any claim this generalizes beyond
the tested generator/mechanisms/perturbation (PAPER\_CONTRACT.md \S2/\S3/\S6).
\end{draftnote}
% EVIDENCE:
% - research/CONTRIBUTION_LOCK.md Sec.6, Sec.11
% - research/PAPER_CONTRACT.md Sec.2, Sec.7
% - research/MANUSCRIPT_ARCHITECTURE.md Sec.2
\end{abstract}

\noindent\textbf{Keywords (candidates only, finalize at E3):} TODO -- e.g.\
classification-mechanism selection; historical label consistency; algorithm
selection; selective classification; representation stability. Do not finalize
without checking each term against research/literature/terminology\_map.md's
adoption list.

% =============================================================================
\section{Introduction}
% =============================================================================
% Structured per the E2 brief's explicit 9-point Introduction plan. Each point
% below is its own \subsection with a draftnote, matching "for every
% subsection include a compact drafting block."

\subsection{1. Real-world motivation}
\begin{draftnote}
PURPOSE: Open with the production observation that prompted this research
question, stated as motivation, never as evidence.\\
QUESTION: Why does this problem matter, concretely?\\
ALLOWED EVIDENCE: METHODOLOGY.md real-vs-synthetic table (the "R3 flip":
production RULES\_FIRST at 91.2\% vs.\ initial synthetic EMBEDDING\_PRIMARY);
research/EVIDENCE\_BASELINE.md \S1 for the exact, caveated production figures.\\
KEY POINTS:
\begin{itemize}
  \item A real production system's architecture-decision rule flipped between
  two single runs (production vs.\ an early synthetic reproduction).
  \item This is stated as an \emph{observation that prompted the question}, not
  as a controlled test of anything.
  \item Every production number used here must carry the "likely understated,
  unverified" caveat from EVIDENCE\_BASELINE.md \S1 (pre-A5-fix production
  figures).
  \item Optional: Table~T1 (production snapshot) may render this as a small,
  visually distinct table -- see the placeholder in Sec.~1.9 area or inline
  here.
\end{itemize}
MUST NOT SAY: that the production case validates, confirms, or independently
tests anything in this paper; that production data is publicly reproducible
(PUBLIC\_RELEASE\_BOUNDARY.md \S3, tier "case-study only / confidential").
\end{draftnote}
% EVIDENCE:
% - METHODOLOGY.md (real-vs-synthetic table)
% - research/EVIDENCE_BASELINE.md Sec.1
% - research/r3_threshold_analysis.md

% TABLE T1 (optional) -- Production case study snapshot.
% Purpose: visually separate "cited" production numbers from "computed here"
% synthetic numbers for the R3-flip observation.
% Source artifact: research/EVIDENCE_BASELINE.md Sec.1 (production),
%   METHODOLOGY.md real-vs-synthetic table (initial synthetic run).
% Expected variables/rows: deterministic-share, weighted ADS, unweighted ADS,
%   R3 decision -- each cell tagged "cited" or "computed here".
% Required or optional: OPTIONAL. Prose may carry this instead; include only
%   if the Introduction reads as needing a visual anchor at E3.

\subsection{2. General problem}
\begin{draftnote}
PURPOSE: State the general pattern this specific observation is an instance
of: architecture/mechanism choice is usually made before, not from, measured
data properties.\\
QUESTION: What is the general problem, independent of this project's specific
domain?\\
ALLOWED EVIDENCE: research/literature/contribution\_status.md (C2 framing);
research/literature/terminology\_map.md.\\
KEY POINTS:
\begin{itemize}
  \item Classification systems typically fix a mechanism class up front.
  \item A measured signal computed from historical data could, in principle,
  inform that choice before any model is built or deployed.
  \item This paper studies whether one specific such signal (historical label
  consistency) actually does so reliably.
\end{itemize}
MUST NOT SAY: that measuring data before choosing an architecture is itself
novel (C2, CHALLENGED -- contribution\_lock.csv row C2).
\end{draftnote}
% EVIDENCE:
% - research/literature/contribution_status.md (C2)
% - research/contribution_lock.csv row C2

\subsection{3. Existing algorithm-selection framing}
\begin{draftnote}
PURPOSE: Situate the general problem against its established name in the
literature, briefly, before Related Work develops it fully.\\
QUESTION: What is the closest existing vocabulary for this general problem?\\
ALLOWED EVIDENCE: research/literature/terminology\_map.md ("Design-time
architecture selection" entry); \citep{rice1976,smithmiles2009,barbudo2023}.\\
KEY POINTS:
\begin{itemize}
  \item Name Rice's Algorithm Selection Problem and its meta-learning /
  workflow-composition descendants as the acknowledged root vocabulary.
  \item State explicitly that this paper operates at the architecture/workflow
  level those literatures describe the field moving toward, not the
  per-instance level Rice originally posed.
  \item Keep this brief -- full positioning is Sec.~2 (Related Work)'s job,
  not the Introduction's.
\end{itemize}
MUST NOT SAY: that this paper's positioning is a new gap beyond what Phase B/C
already found (research/literature/contribution\_status.md).
\end{draftnote}
% EVIDENCE:
% - research/literature/terminology_map.md
% - research/literature/prior_art_map.md

\subsection{4. Specific research question}
\begin{draftnote}
PURPOSE: State H1's parent research question exactly, in the paper's own
words, before any experimental detail.\\
QUESTION: What, precisely, does this paper ask?\\
ALLOWED EVIDENCE: research/CONTRIBUTION\_LOCK.md \S11.A (verbatim question).\\
KEY POINTS:
\begin{itemize}
  \item "Can historical decision consistency, measured before deployment, be
  used to select between qualitatively different classification mechanisms
  (exact-match rules vs.\ retrieval-based fuzzy matching)?"
  \item State the four preconditions explicitly here, not deferred to
  Limitations: a history of repeated decisions, observable labels, measurable
  consistency, sufficient historical coverage.
\end{itemize}
MUST NOT SAY: a broader question than this (e.g.\ "...select any AI system
architecture"); PAPER\_CONTRACT.md \S6 (generalization rule).
\end{draftnote}
% EVIDENCE:
% - research/CONTRIBUTION_LOCK.md Sec.11.A
% - STATE.md ("Literature review + paper-positioning conclusions")

\subsection{5. Why historical consistency is attractive as a signal}
\begin{draftnote}
PURPOSE: Give the reader the intuitive case for why this signal seemed
promising, before showing where it breaks down.\\
QUESTION: Why would anyone expect this signal to work?\\
ALLOWED EVIDENCE: TECHNICAL\_REPORT.md \S2.2 (ADS motivation prose, used only
for framing, not for numbers -- per PAPER\_CONTRACT.md \S4, verified artifacts
outrank TECHNICAL\_REPORT.md prose wherever they might differ).\\
KEY POINTS:
\begin{itemize}
  \item If a product has almost always been booked to one account historically,
  that looks like strong evidence the mapping is "solved" and a simple
  mechanism (exact lookup) should suffice.
  \item The natural hypothesis: high consistency $\Rightarrow$ a simple,
  precedent-based mechanism should be preferred over a more flexible one.
  \item This is exactly the intuition this paper tests, and partly falsifies.
\end{itemize}
MUST NOT SAY: any claim that this intuition is confirmed before Results is
introduced.
\end{draftnote}
% EVIDENCE:
% - TECHNICAL_REPORT.md Sec.2.2 (framing only, not numbers)
% - research/CONTRIBUTION_LOCK.md Sec.2 step 3

\subsection{6. What is actually tested}
\begin{draftnote}
PURPOSE: Compress the experimental design into two or three sentences so a
reader has the shape of the test before the findings.\\
QUESTION: What, concretely, does the experiment compare?\\
ALLOWED EVIDENCE: research/EXPERIMENT\_1\_REDESIGN\_REVIEW.md \S2, \S6.\\
KEY POINTS:
\begin{itemize}
  \item A pre-registered, 240-condition synthetic factorial experiment.
  \item Two isolated mechanisms: exact-match rules vs.\ rapidfuzz-based
  retrieval (cutoff=75).
  \item Two lexical conditions (CLEAN / VARIED) crossed with 6 nominal
  consistency targets and 20 seeds.
  \item The LLM mechanism is excluded, for a stated, principled reason (not an
  oversight) -- full justification in Sec.~4.
\end{itemize}
MUST NOT SAY: any result yet -- this subsection previews scope, not findings.
\end{draftnote}
% EVIDENCE:
% - research/EXPERIMENT_1_REDESIGN_REVIEW.md Sec.2, Sec.6, Sec.10

\subsection{7. Main findings}
\begin{draftnote}
PURPOSE: State the two-part finding (6a/6b) plainly, exactly once, before the
paper's Contribution statement makes it formal.\\
QUESTION: What did we find, in one paragraph?\\
ALLOWED EVIDENCE: research/CONTRIBUTION\_LOCK.md \S6.\\
KEY POINTS:
\begin{itemize}
  \item 6a: realized ADS strongly predicts each mechanism's own accuracy.
  \item 6b: ADS does not predict which mechanism wins; the winner is a
  constant function of the (manipulated) lexical condition.
  \item These two claims are never merged into one sentence.
\end{itemize}
MUST NOT SAY: "ADS predicts mechanism suitability" (unqualified) -- the single
highest-risk collapse this entire manuscript exists to avoid
(PAPER\_CONTRACT.md \S3 row 13).
\end{draftnote}
% EVIDENCE:
% - research/CONTRIBUTION_LOCK.md Sec.4, Sec.6

\subsection{8. Contribution statement}
\begin{draftnote}
PURPOSE: The formal, citable contribution sentence(s) -- must match
CONTRIBUTION\_LOCK.md \S6 exactly, no stronger.\\
QUESTION: What, precisely, does this paper contribute?\\
ALLOWED EVIDENCE: research/CONTRIBUTION\_LOCK.md \S6 (6a, 6b, Synthesis),
verbatim in spirit.\\
KEY POINTS:
\begin{itemize}
  \item Reuse CONTRIBUTION\_LOCK.md \S6's exact wording (6a, 6b, Synthesis) --
  do not paraphrase into something stronger.
  \item State plainly that this is not a new metric, not a new architecture,
  and not a validated selection method (PAPER\_CONTRACT.md \S8).
\end{itemize}
MUST NOT SAY: any of PAPER\_CONTRACT.md \S3's 16 forbidden claims -- this
subsection is the single highest-scrutiny paragraph in the entire paper and
should be checked against that list line by line at E4.
\end{draftnote}
% EVIDENCE:
% - research/CONTRIBUTION_LOCK.md Sec.6 (exact wording to reuse)
% - research/PAPER_CONTRACT.md Sec.8

\subsection{9. Paper roadmap}
\begin{draftnote}
PURPOSE: One short paragraph orienting the reader to the remaining eight
sections.\\
QUESTION: Where does the paper go from here?\\
ALLOWED EVIDENCE: This document's own section list (Sec.~2--9 below).\\
KEY POINTS:
\begin{itemize}
  \item Related Work $\to$ Problem Setting $\to$ Experimental Design $\to$
  Results $\to$ Discussion $\to$ Limitations $\to$ Future Work $\to$
  Conclusion.
\end{itemize}
MUST NOT SAY: nothing to forbid here -- this is pure orientation, not a claim.
\end{draftnote}


% =============================================================================
\section{Related Work}
% =============================================================================
% Structured around the exact research families research/literature/
% contribution_status.md, prior_art_map.md, and terminology_map.md already
% established -- no new family invented, no new gap claimed. Citation keys
% only; no new literature search performed (per the E2 brief).

\subsection{Cluster Purity and Majority-Vote Agreement}
\begin{draftnote}
PRIOR ART: \citep{manning2008,amigo2009} define cluster purity as
$\mathrm{purity}(\mathrm{cluster}_k) = \max_j(\text{count of class }j\text{ in
cluster }k) / |\mathrm{cluster}_k|$ -- the identical closed-form expression as
ADS under "cluster" $\leftrightarrow$ "item's historical booking multiset".
\citep{dawidskene1979} and its descendant crowdsourcing/truth-inference
literature use the same raw majority-vote-agreement proportion as the
universal naive baseline since 1979.\\
OUR NON-CLAIM: that ADS is a novel metric. This is a closed question, not an
open one (contribution\_lock.csv row C1, CHALLENGED).\\
NARROW DELTA: none at the metric level. Any delta this paper has is entirely
in what the metric is \emph{used for} (mechanism selection), covered in the
next subsection, not in the formula itself.
\end{draftnote}
% EVIDENCE:
% - research/literature/ads_metric_prior_art.md
% - research/literature/ads_metric_prior_art.csv rows G1-01..G1-04
% - research/CONTRIBUTION_LOCK.md Sec.3 (C1)

\subsection{Algorithm Selection and Meta-Learning}
\begin{draftnote}
PRIOR ART: \citep{rice1976} originates the Algorithm Selection Problem --
select a single algorithm per problem instance. \citep{smithmiles2009}
unifies this with meta-learning; later work (Ali \& Smith 2006; Khan et al.\
2020, ledger rows B2-10/B2-11) selects a generic ML algorithm from historical
evidence in the same application domain (classification).\\
OUR NON-CLAIM: that design-time selection from historical evidence is itself
new (C2, CHALLENGED, contribution\_lock.csv row C2).\\
NARROW DELTA: this paper studies whether a \emph{label-consistency-specific}
signal (not generic meta-features, not performance search) predicts a
\emph{qualitative mechanism-class} choice -- the C2b framing -- and reports
exactly where that narrower claim holds and where it fails.
\end{draftnote}
% EVIDENCE:
% - research/literature/citation_ledger.csv rows B1-01..B1-05, B2-01, B2-10, B2-11
% - research/literature/contribution_status.md (C2, C2b)

\subsection{AutoML and Workflow Composition}
\begin{draftnote}
PRIOR ART: \citep{barbudo2023} documents the field's shift from single-
algorithm selection to automated, benchmarked-performance-driven workflow
composition -- "the algorithm selection problem has gradually been superseded
by the challenge of workflow composition," the single closest terminological
echo of this paper's own framing found in the literature sweep.\\
OUR NON-CLAIM: that this paper's simple, interpretable threshold rule (R3) is
a competing AutoML method.\\
NARROW DELTA: evidence type (label-consistency vs.\ benchmarked search) and
interpretability (a named threshold vs.\ an optimization search) -- stated as
a difference in kind, not a claim of superiority.
\end{draftnote}
% EVIDENCE:
% - research/literature/citation_ledger.csv row B2-02
% - research/literature/prior_art_map.md (Tier 2, "Design-time architecture selection")

\subsection{Self-Designed and Instance-Optimized Data Systems}
\begin{draftnote}
PRIOR ART: \citep{idreoskraska2019} is the strongest non-ML structural
analog: measure historical evidence once, choose system composition before
serving -- but the evidence type is workload/query \emph{frequency}, and the
paper's own thesis is continuous self-adaptation, explicitly rejecting a
one-time human design gate (the opposite temporal structure).\\
OUR NON-CLAIM: that a one-shot, pre-deployment design gate is itself a
contribution (already the opposite of this literature's stated thesis, cited
as contrast, not as an unsolved problem this paper solves).\\
NARROW DELTA: this paper's experiment is about the specific signal type
(label consistency, not workload frequency) and its specific limitation
(blindness to representation instability), not a defense of the one-shot-gate
design pattern generally.
\end{draftnote}
% EVIDENCE:
% - research/literature/citation_ledger.csv rows B7-01, B7-05, B7-06
% - research/literature/terminology_map.md ("Design-time architecture selection")

\subsection{Reject Option and Selective Classification}
\begin{draftnote}
PRIOR ART: \citep{chow1970} originates the single-threshold accept/reject
decision on live posterior confidence; \citep{elyaniv2010} formalizes the
risk-coverage tradeoff; \citep{hendrickx2024} is the most comprehensive recent
survey of the entire field.\\
OUR NON-CLAIM: that Experiment 1 introduces a new reject-option variant.\\
NARROW DELTA: not directly relevant to Experiment 1 at all (Experiment 1 tests
a design-time selection rule, not a runtime reject threshold) -- mentioned
only to establish the design-time vs.\ inference-time distinction this paper's
scope sits on one side of.
\end{draftnote}
% EVIDENCE:
% - research/literature/citation_ledger.csv rows B4-01, B4-02, B4-05
% - STATE.md (settled design-time/inference-time positioning)

\subsection{Learning to Defer}
\begin{draftnote}
PRIOR ART: \citep{mozannarsontag2020} is the foundational modern
learning-to-defer paper: a trained/rule-based function decides whether to
defer to a human, informed by historical decision data, as a single binary
per-item decision.\\
OUR NON-CLAIM: that this experiment's mechanism-selection rule is a
learning-to-defer instance.\\
NARROW DELTA: establishes the same design-time/inference-time boundary this
paper's scope sits on one side of; the (untested here) production cascade
sits closer to this literature, but that system is out of scope for
Experiment 1.
\end{draftnote}
% EVIDENCE:
% - research/literature/citation_ledger.csv rows B5-01, B5-02
% - research/literature/terminology_map.md ("Runtime confidence cascade")

\subsection{Model Cascades and LLM Routing}
\begin{draftnote}
PRIOR ART: \citep{frugalgpt2023} is the closest structural match to a
multi-tier, historically-calibrated cascade; \citep{rankgpt2023} establishes
LLM-reranks-a-pre-fetched-candidate-list as commodity IR technique since 2023.\\
OUR NON-CLAIM: that Experiment 1's two-mechanism comparison is a cascade
contribution.\\
NARROW DELTA: not directly relevant -- the LLM mechanism is explicitly
excluded from Experiment 1 (Sec.~4); this family is cited only for
completeness of the design-time/runtime distinction, not engaged as a
comparison point for the actual experiment reported here.
\end{draftnote}
% EVIDENCE:
% - research/literature/citation_ledger.csv rows B3-02, B6-02
% - research/EXPERIMENT_1_REDESIGN_REVIEW.md Sec.10 (LLM exclusion rationale)

\subsection{Domain-Specific Practice}
\begin{draftnote}
PRIOR ART: \citep{jorgensenigel2021} empirically demonstrates, in the same
application domain, that a global classifier generalizes far worse across
companies than per-company models -- the exact phenomenon the production
system's cross-company-consistency threshold encodes. Industry sources (Ken
From Finance, Peakflo, Ramp -- ledger rows B8-04/05/06) independently describe
similar informal historical-consistency-audit practice in production AP
automation.\\
OUR NON-CLAIM: that the application domain itself is unprecedented, or that no
vendor measures before choosing (industry sources directly contradict that
framing -- see the known, still-open TECHNICAL\_REPORT.md \S5 correction item
tracked in research/MANUSCRIPT\_CLAIM\_EVIDENCE\_MAP.md \S4).\\
NARROW DELTA: the domain remains academically under-served (no peer-reviewed
SAF-T/D406-specific ML paper found) -- a narrow, honest niche-novelty claim,
much weaker than any methodological novelty claim, and not the paper's actual
scientific contribution.
\end{draftnote}
% EVIDENCE:
% - research/literature/citation_ledger.csv rows B8-01, B8-04, B8-05, B8-06
% - research/literature/contribution_status.md (C8)
% - research/MANUSCRIPT_CLAIM_EVIDENCE_MAP.md Sec.4

% TABLE T2 -- Literature/contribution positioning.
% Purpose: one-glance summary of every Related Work subsection's
%   PRIOR ART / OUR NON-CLAIM / NARROW DELTA columns.
% Source artifact: research/literature/contribution_status.md,
%   research/literature/prior_art_map.md, research/literature/terminology_map.md.
% Expected variables/rows: one row per research family (8 rows, matching the
%   8 subsections above).
% Required or optional: REQUIRED (per research/MANUSCRIPT_ARCHITECTURE.md
%   Sec.8) -- reuses exactly the 8-row table already drafted there Sec.10;
%   do not re-derive at E3, transcribe it.


% =============================================================================
\section{Problem Setting and Signal Definition}
% =============================================================================
% Merged Problem Setting + Method per research/MANUSCRIPT_ARCHITECTURE.md
% Sec.3.1's explicit merge decision: ADS's formal definition IS the formal
% problem setting here; there is no separate "method" this paper contributes
% (ADS is explicitly not a novel-metric claim).

\subsection{Historical Product-Level Decisions and Account Labels}
\begin{draftnote}
PURPOSE: Define the data object every downstream signal and mechanism
operates on, in general (domain-agnostic) terms.\\
QUESTION: What is a "historical decision" in this paper's sense?\\
ALLOWED EVIDENCE: TECHNICAL\_REPORT.md \S2.1 (framing only); research/
EXPERIMENT\_1\_DATA\_DICTIONARY.md ("Unit of observation").\\
KEY POINTS:
\begin{itemize}
  \item A repeated, historically observed mapping from an item (a "product")
  to a label (an "account"), with an occurrence count.
  \item This paper's method is domain-agnostic: nothing below references
  Romania, fiscal documents, or invoices specifically.
  \item State the four preconditions again, briefly, in formal terms (repeated
  decisions, observable labels, measurable consistency, sufficient coverage).
\end{itemize}
MUST NOT SAY: anything domain-specific here -- that framing belongs only in
the Introduction's motivation, not in the formal problem setting.
\end{draftnote}
% EVIDENCE:
% - research/EXPERIMENT_1_DATA_DICTIONARY.md ("Unit of observation")
% - STATE.md (four preconditions)

\subsection{The Automated Determinism Score (ADS)}
\begin{draftnote}
PURPOSE: Give the exact formula, paired immediately with its equivalence
citation in the same paragraph -- never a bare definition.\\
QUESTION: What, precisely, is the signal this paper studies?\\
ALLOWED EVIDENCE: TECHNICAL\_REPORT.md \S2.2 (formula only);
\citep{manning2008,amigo2009,dawidskene1979}.\\
KEY POINTS:
\begin{itemize}
  \item State Equation~\ref{eq:ads} and its two aggregates (weighted,
  unweighted) if needed for later sections.
  \item Immediately follow the definition with: "this is mathematically
  identical to cluster purity \citep{manning2008,amigo2009} and the raw
  majority-vote-agreement baseline \citep{dawidskene1979}" -- per
  terminology\_map.md's explicit warning that silence here is "the single
  highest-risk omission a literature-aware reviewer would flag."
  \item ADS is descriptive, not novel, and not this paper's contribution.
\end{itemize}
MUST NOT SAY: "novel metric," "our metric," or any bare ADS definition
without the equivalence sentence in the same paragraph.
\end{draftnote}
% EVIDENCE:
% - TECHNICAL_REPORT.md Sec.2.2 (formula)
% - research/literature/ads_metric_prior_art.md
% - research/PAPER_CONTRACT.md Sec.3 row 1

\begin{equation}
\label{eq:ads}
\mathrm{ADS}(p) \;=\; \frac{\max_i c_i}{\sum_i c_i}
\end{equation}
% EQUATION E1 -- ADS definition. $c_1,\dots,c_k$ are the historical occurrence
% counts of item $p$ across its $k$ observed accounts. Not presented as a
% novel mathematical contribution (see draftnote above and
% research/MANUSCRIPT_ARCHITECTURE.md Sec.9, row E1).

\subsection{Exact-Match Rules Mechanism}
\begin{draftnote}
PURPOSE: Define the first of the two compared mechanisms precisely enough
that Sec.~4's experimental description needs no further clarification.\\
QUESTION: What, exactly, does "rules" mean in this experiment?\\
ALLOWED EVIDENCE: research/EXPERIMENT\_1\_REDESIGN\_REVIEW.md \S11;
research/EXPERIMENT\_1\_DATA\_DICTIONARY.md (\texttt{rules\_whole\_set\_accuracy},
\texttt{rules\_coverage}).\\
KEY POINTS:
\begin{itemize}
  \item Exact company-scoped string lookup, then exact global lookup; any hit
  is the mechanism's answer; no hit = abstain.
  \item No confidence threshold to calibrate -- binary exact-match/no-match.
  \item Abstention counts as incorrect under the whole-set-accuracy metric
  (Sec.~4) -- no partial credit.
  \item Coverage is always $<1.0$ across all 240 conditions (a real, reported
  asymmetry with retrieval, not normalized away).
\end{itemize}
MUST NOT SAY: that this mechanism is the shipped production cascade's Tier 1
(it is an isolated, simplified test harness mechanism, not the multi-tier
cascade -- PAPER\_CONTRACT.md \S6).
\end{draftnote}
% EVIDENCE:
% - research/EXPERIMENT_1_REDESIGN_REVIEW.md Sec.11
% - research/EXPERIMENT_1_DATA_DICTIONARY.md (rules_* columns)

\subsection{Fuzzy/Similarity Retrieval Mechanism}
\begin{draftnote}
PURPOSE: Define the second mechanism, including why it is named
"retrieval," not "embedding," per the redesign review's explicit renaming
decision.\\
QUESTION: What, exactly, does "retrieval" mean in this experiment?\\
ALLOWED EVIDENCE: research/EXPERIMENT\_1\_REDESIGN\_REVIEW.md \S9, \S11;
research/EXPERIMENT\_1\_DATA\_DICTIONARY.md (\texttt{retrieval\_*} columns).\\
KEY POINTS:
\begin{itemize}
  \item rapidfuzz lexical-similarity match (\texttt{WRatio}), company-scoped
  then global, used as the \emph{primary} classifier -- not a fallback gated
  behind a rules-miss.
  \item Cutoff = 75, calibrated once (Sec.~4), frozen for the entire run.
  \item Coverage is $1.0$ in all 240 conditions -- retrieval never abstains at
  this cutoff.
  \item Explicitly \emph{not} a semantic embedding model -- state the renaming
  decision (redesign review \S9) so a reader does not conflate this with a
  vector-embedding retrieval system.
\end{itemize}
MUST NOT SAY: "embedding-primary" or any language implying a trained
embedding model was used (explicitly rejected terminology, redesign review
\S9).
\end{draftnote}
% EVIDENCE:
% - research/EXPERIMENT_1_REDESIGN_REVIEW.md Sec.9, Sec.11
% - research/EXPERIMENT_1_DATA_DICTIONARY.md (retrieval_* columns)

\subsection{Mechanism Accuracy vs. Mechanism Ranking}
\begin{draftnote}
PURPOSE: State the central distinction this entire paper is organized around,
before any result is shown, so Results (Sec.~5) can simply invoke it.\\
QUESTION: What is the difference between "ADS predicts accuracy" and "ADS
predicts ranking," precisely?\\
ALLOWED EVIDENCE: research/CONTRIBUTION\_LOCK.md \S4.\\
KEY POINTS:
\begin{itemize}
  \item "ADS predicts mechanism accuracy": a correlational claim about each
  mechanism's own performance level (Sec.~5.2).
  \item "ADS predicts mechanism ranking": a claim about \emph{which}
  mechanism wins (Sec.~5.3) -- a different question with a different (here,
  negative) answer.
  \item These are not the same claim, do not stand or fall together, and this
  paper never merges them (CONTRIBUTION\_LOCK.md \S4, verbatim requirement).
\end{itemize}
MUST NOT SAY: any sentence that could be read as treating these as
interchangeable.
\end{draftnote}
% EVIDENCE:
% - research/CONTRIBUTION_LOCK.md Sec.4
% - research/PAPER_CONTRACT.md Sec.2 rows 3-4


% =============================================================================
\section{Experimental Design}
% =============================================================================

\subsection{Research Hypothesis}
\begin{draftnote}
PURPOSE: State H1 (revised) exactly as pre-registered, before any design
detail.\\
QUESTION: What, precisely, was hypothesized, and when?\\
ALLOWED EVIDENCE: research/EXPERIMENT\_1\_REDESIGN\_REVIEW.md \S2.\\
KEY POINTS:
\begin{itemize}
  \item H1 (revised): under controlled synthetic conditions, higher realized
  historical decision consistency will be associated with predictable changes
  in relative mechanism performance, such that a pre-specified, frozen
  consistency-based decision rule (R3) will agree with the empirically
  best-performing mechanism more often than chance.
  \item State plainly this was pre-registered \emph{before} any data existed,
  with falsification criteria fixed in advance (Sec.~4.13).
\end{itemize}
MUST NOT SAY: a stronger hypothesis than this (e.g.\ "monotonic dominance in
every band") -- H1 was deliberately weaker than an earlier draft phrasing, per
the redesign review's own explicit rationale.
\end{draftnote}
% EVIDENCE:
% - research/EXPERIMENT_1_REDESIGN_REVIEW.md Sec.2

\subsection{Synthetic Generator}
\begin{draftnote}
PURPOSE: Describe the generator at the scale and scope actually used.\\
QUESTION: What generates the data every condition is built from?\\
ALLOWED EVIDENCE: research/EXPERIMENT\_1\_REDESIGN\_REVIEW.md \S4, \S6.\\
KEY POINTS:
\begin{itemize}
  \item 60 companies, 1,200-product vocabulary, fixed across every condition.
  \item Per-call \texttt{random.Random(seed)} isolation (not a single mutated
  global RNG) -- a documented, fixed implementation requirement, since a
  multi-seed sweep needs independent, reproducible streams per seed.
  \item Cross-company alignment fixed at 0.695 (an existing, evidence-derived
  production value) -- not swept, to avoid a second independent variable.
\end{itemize}
MUST NOT SAY: that this generator is calibrated against real invoice/OCR
error data (it is not -- Sec.~7.2).
\end{draftnote}
% EVIDENCE:
% - research/EXPERIMENT_1_REDESIGN_REVIEW.md Sec.4, Sec.5

\subsection{Data-Generating Factors}
\begin{draftnote}
PURPOSE: Enumerate the two controlled factors of the actual factorial design.\\
QUESTION: What varies systematically across conditions?\\
ALLOWED EVIDENCE: research/EXPERIMENT\_1\_FINAL\_RESULTS.md \S3.\\
KEY POINTS:
\begin{itemize}
  \item Factor 1 -- nominal target \texttt{deterministic\_share}: 6 levels
  actually run, $\{0.00, 0.20, 0.30, 0.50, 0.75, 1.00\}$ (the frozen final run
  -- note this differs from the wider calibration-only grid used for Gate 1,
  Sec.~4.12; do not conflate the two).
  \item Factor 2 -- lexical condition: CLEAN vs.\ VARIED, 2 levels.
  \item 6 $\times$ 2 $\times$ 20 seeds = 240 conditions, exactly.
\end{itemize}
MUST NOT SAY: the calibration grid's 17-point target range (Sec.~4.12) as if
it were the frozen run's design -- the two are different, both real,
never-conflated artifacts.
\end{draftnote}
% EVIDENCE:
% - research/EXPERIMENT_1_FINAL_RESULTS.md Sec.3
% - research/EXPERIMENT_1_EVIDENCE_CHECKPOINT.md Sec.3

\subsection{Realized ADS vs. Target Deterministic Share}
\begin{draftnote}
PURPOSE: Explicitly distinguish the generator's input knob from the measured
outcome -- per the E2 brief, this distinction must never be collapsed
anywhere in the manuscript.\\
QUESTION: What is the actual independent variable, and how does it differ
from the generator's control parameter?\\
ALLOWED EVIDENCE: research/EXPERIMENT\_1\_REDESIGN\_REVIEW.md \S3, \S7;
research/EXPERIMENT\_1\_DATA\_DICTIONARY.md (\texttt{realized\_det\_pct} vs.\
\texttt{target\_deterministic\_share}).\\
KEY POINTS:
\begin{itemize}
  \item \texttt{target\_deterministic\_share} is the generator's \emph{input
  knob} -- a generation-time control parameter, never itself an axis of any
  primary result.
  \item \texttt{realized\_det\_pct} ("realized ADS") is the \emph{measured,
  train-only} outcome, grouped by the stable \texttt{product\_code}, computed
  \emph{after} generation. This is the actual primary independent variable
  and the only quantity R3 (Sec.~4.10) thresholds.
  \item The two usually but do not always coincide: at target=0.50, individual
  seeds' realized ADS straddles the 0.70 R3 boundary (range 0.659--0.727
  across 20 seeds), so roughly half the seeds sharing one nominal target land
  in each of two different R3 bands.
  \item Observed realized-ADS range across the frozen run: 0.44--0.93,
  structurally capped below $\sim$0.91 by the fixed cross-company-alignment
  nuisance parameter (Sec.~4.2) -- the "deep rules-first" region ($\geq$0.93)
  was never reachable.
  \item Realized ADS is invariant to the lexical condition by construction
  (grouped by \texttt{product\_code}, never the perturbable surface string) --
  confirmed exactly in the data (byte-identical means/stds between CLEAN and
  VARIED at every target).
\end{itemize}
MUST NOT SAY: any result, table, or figure axis that uses
\texttt{target\_deterministic\_share} where \texttt{realized\_det\_pct} is the
scientifically correct axis (Sec.~5 uses realized ADS exclusively).
\end{draftnote}
% EVIDENCE:
% - research/EXPERIMENT_1_REDESIGN_REVIEW.md Sec.3, Sec.7
% - research/EXPERIMENT_1_DATA_DICTIONARY.md ("What realized ADS is, precisely")
% - research/EXPERIMENT_1_POSTHOC_ANALYSIS.md Sec.3

\subsection{The CLEAN Condition}
\begin{draftnote}
PURPOSE: Define the no-noise baseline condition.\\
QUESTION: What does "CLEAN" mean, exactly?\\
ALLOWED EVIDENCE: research/EXPERIMENT\_1\_REDESIGN\_REVIEW.md \S8.\\
KEY POINTS:
\begin{itemize}
  \item The existing generator's default output, byte-identical to today's
  committed baseline -- no surface-string perturbation applied.
  \item Under CLEAN, rules and retrieval are statistically tied (within
  $\delta=0.02$) at every one of the 6 targets -- zero conditions produce a
  defined empirical winner (Sec.~5.3).
\end{itemize}
MUST NOT SAY: that CLEAN "confirms" mechanism equivalence in general -- only
for this generator, absent lexical noise (PAPER\_CONTRACT.md \S3 row 10).
\end{draftnote}
% EVIDENCE:
% - research/EXPERIMENT_1_REDESIGN_REVIEW.md Sec.8
% - research/EXPERIMENT_1_FINAL_RESULTS.md Sec.4 (CLEAN rows)

\subsection{The VARIED Condition and Lexical Perturbation}
\begin{draftnote}
PURPOSE: Define the noise condition and the exact, fixed transform algorithm.\\
QUESTION: What, exactly, is perturbed, how, and how often?\\
ALLOWED EVIDENCE: research/EXPERIMENT\_1\_REDESIGN\_REVIEW.md \S8.\\
KEY POINTS:
\begin{itemize}
  \item Applied only to the display/lookup string (\texttt{normalized\_
  product}); the ground-truth key (\texttt{product\_code}) is untouched --
  semantic identity is preserved by construction, not by convention.
  \item Five fixed transform types: case variation, punctuation variation,
  token reorder, abbreviation, whitespace variation.
  \item \texttt{P\_TRANSFORM=0.3} -- an empirically pilot-tuned value (not a
  first-principles constant), frozen before the run, never re-tuned after
  seeing mechanism-accuracy results.
  \item Fully reproducible: deterministic per-line sub-seed derivation, no
  extra state to store.
\end{itemize}
MUST NOT SAY: that \texttt{P\_TRANSFORM=0.3} represents measured real-world
OCR/typo error rates (PAPER\_CONTRACT.md \S3 row 11) -- it is a pilot-tuned
synthetic stand-in only.
\end{draftnote}
% EVIDENCE:
% - research/EXPERIMENT_1_REDESIGN_REVIEW.md Sec.8
% - research/EXPERIMENT_1_CALIBRATION_REPORT.md (P_TRANSFORM calibration)

\subsection{Mechanisms and Retrieval Cutoff}
\begin{draftnote}
PURPOSE: Cross-reference Sec.~3's mechanism definitions with the one
calibrated parameter (retrieval cutoff) and state how it was fixed.\\
QUESTION: How was the retrieval cutoff chosen, and is it re-tuned anywhere?\\
ALLOWED EVIDENCE: research/EXPERIMENT\_1\_REDESIGN\_REVIEW.md \S12
(calibration protocol); research/EXPERIMENT\_1\_CALIBRATION\_REPORT.md.\\
KEY POINTS:
\begin{itemize}
  \item Only \texttt{retrieval\_only} has a tunable threshold; \texttt{rules\_
  only} has none.
  \item Calibrated once, on a dedicated, separate generation run (its own
  seed, distinct from every frozen-sweep seed), before the frozen sweep.
  \item Selected cutoff: 75 (candidate range $\{70,75,80,85,90\}$, chosen by
  maximizing whole-set accuracy subject to a 30\% coverage floor).
  \item Frozen for the entire run -- never re-tuned per band, per lexical
  condition, or after seeing results.
\end{itemize}
MUST NOT SAY: that the cutoff was calibrated on the same data the final
comparison ran on (it was not -- a separate dedicated calibration run).
\end{draftnote}
% EVIDENCE:
% - research/EXPERIMENT_1_REDESIGN_REVIEW.md Sec.12
% - research/EXPERIMENT_1_CALIBRATION_REPORT.md

\subsection{Seed Structure}
\begin{draftnote}
PURPOSE: State the exact seed manifest.\\
QUESTION: How many seeds, and how were they chosen?\\
ALLOWED EVIDENCE: research/EXPERIMENT\_1\_FINAL\_RESULTS.md \S3;
research/EXPERIMENT\_1\_EVIDENCE\_CHECKPOINT.md \S3.\\
KEY POINTS:
\begin{itemize}
  \item 20 seeds, $\{31001, \ldots, 31020\}$, identical set across every one
  of the 12 (target, lexical) cells.
  \item A documented, deterministic range rule, never hand-picked after seeing
  any result.
\end{itemize}
MUST NOT SAY: nothing to forbid -- purely a configuration fact.
\end{draftnote}
% EVIDENCE:
% - research/EXPERIMENT_1_FINAL_RESULTS.md Sec.3
% - research/EXPERIMENT_1_EVIDENCE_CHECKPOINT.md Sec.3

\subsection{Train/Test Separation}
\begin{draftnote}
PURPOSE: State the leakage-safety protocol precisely -- this is the
credibility backbone of the entire realized-ADS measurement.\\
QUESTION: What data does the consistency signal see, and what does the
evaluation see?\\
ALLOWED EVIDENCE: research/EXPERIMENT\_1\_REDESIGN\_REVIEW.md \S13;
research/EXPERIMENT\_1\_EVIDENCE\_CHECKPOINT.md \S11.\\
KEY POINTS:
\begin{itemize}
  \item \texttt{realized\_det\_pct} is computed exclusively from the train
  split -- mirroring real deployment temporal logic (a design-time decision
  can only see history, never the held-out future).
  \item The mechanism knowledge base is built only from train-split rows.
  \item Every accuracy/coverage number reported is computed exclusively on
  that condition's test rows.
  \item Verified in code, not merely asserted: \texttt{consistency.py} filters
  to \texttt{split\_of(r)=="train"} before any aggregation; test account
  labels never reach the mechanism-selection rule.
\end{itemize}
MUST NOT SAY: nothing to forbid -- this subsection is purely a leakage-safety
guarantee statement.
\end{draftnote}
% EVIDENCE:
% - research/EXPERIMENT_1_REDESIGN_REVIEW.md Sec.13
% - research/EXPERIMENT_1_EVIDENCE_CHECKPOINT.md Sec.11 (integrity checks)

\subsection{Paired Comparison and Winner/Tie Definition}
\begin{draftnote}
PURPOSE: Define exactly how a condition's "empirical winner" is determined --
a defined statistical procedure, not a judgment call.\\
QUESTION: How is "which mechanism won" decided, precisely?\\
ALLOWED EVIDENCE: research/EXPERIMENT\_1\_REDESIGN\_REVIEW.md \S14;
research/EXPERIMENT\_1\_DATA\_DICTIONARY.md (\texttt{empirical\_winner}).\\
KEY POINTS:
\begin{itemize}
  \item Both mechanisms are scored on \emph{identical} held-out test lines
  within a condition -- paired, not independent, observations.
  \item Winner: strictly higher whole-set accuracy, \emph{unless} the paired
  bootstrap 95\% CI on the accuracy difference overlaps zero within
  $\delta=0.02$, in which case the condition is scored a TIE (a distinct
  category, never broken toward either mechanism).
  \item State Equation~\ref{eq:winner} here.
\end{itemize}
MUST NOT SAY: that a tie is silently dropped or force-assigned to one
mechanism.
\end{draftnote}
% EVIDENCE:
% - research/EXPERIMENT_1_REDESIGN_REVIEW.md Sec.14
% - research/EXPERIMENT_1_DATA_DICTIONARY.md (empirical_winner column)

\begin{equation}
\label{eq:winner}
\mathrm{winner} =
\begin{cases}
\text{rules} & \text{if } \mathrm{CI}_{95\%}(\mathrm{acc}_{\mathrm{rules}} - \mathrm{acc}_{\mathrm{retrieval}}) \text{ entirely} > +\delta \\
\text{retrieval} & \text{if } \mathrm{CI}_{95\%}(\mathrm{acc}_{\mathrm{rules}} - \mathrm{acc}_{\mathrm{retrieval}}) \text{ entirely} < -\delta \\
\text{tie} & \text{otherwise}
\end{cases}
\end{equation}
% EQUATION E3 -- empirical winner / tie definition. delta = 0.02 (frozen,
% Gate 3). Not a methodological contribution -- standard practical-equivalence
% testing (research/MANUSCRIPT_ARCHITECTURE.md Sec.9, row E3).

\subsection{Statistical Analysis}
\begin{draftnote}
PURPOSE: State the primary metric, the aggregation rule, and the chance
baseline used for the binomial test.\\
QUESTION: How is agreement between R3 and the empirical winner measured and
tested?\\
ALLOWED EVIDENCE: research/EXPERIMENT\_1\_REDESIGN\_REVIEW.md \S14;
research/EXPERIMENT\_1\_EVIDENCE\_CHECKPOINT.md \S6.\\
KEY POINTS:
\begin{itemize}
  \item Primary metric: whole-set accuracy (abstention counts as incorrect) --
  a single, pre-specified number, not an ad hoc composite.
  \item Agreement rate = (\# conditions where R3-selected mechanism ==
  empirical winner) / (\# conditions with a \emph{defined} comparison,
  excluding TIE and N/A rows).
  \item Chance baseline is \textbf{0.5}, not 1/3 -- the LLM leg is excluded
  from the primary comparison, so R3 discriminates between exactly two
  mechanisms wherever a comparison is defined.
  \item Test statistic: one-sided/two-sided exact binomial test against
  $p=0.5$, plus a Wilson CI on the observed agreement rate.
\end{itemize}
MUST NOT SAY: a 1/3 chance baseline (would understate how hard the rule has
to work to look good -- explicitly flagged as a framing error to avoid in the
redesign review \S14).
\end{draftnote}
% EVIDENCE:
% - research/EXPERIMENT_1_REDESIGN_REVIEW.md Sec.14
% - research/EXPERIMENT_1_EVIDENCE_CHECKPOINT.md Sec.6

\subsection{Calibration}
\begin{draftnote}
PURPOSE: Summarize the three pre-registered calibration gates closed before
the frozen run.\\
QUESTION: What was calibrated, and how, before any frozen data was generated?\\
ALLOWED EVIDENCE: research/EXPERIMENT\_1\_CALIBRATION\_REPORT.md.\\
KEY POINTS:
\begin{itemize}
  \item Gate 1 -- ADS calibration: a dense 17-point target grid, 10 seeds per
  target, mechanism-blind (no \texttt{mechanisms.py} import at all in the
  calibration script).
  \item Gate 2 -- retrieval-cutoff calibration (Sec.~4.7).
  \item Gate 3 -- frozen empirical-winner/tie/CI definition (Sec.~4.10).
  \item State explicitly that the Gate-1 calibration grid (17 points, 0.00 to
  1.00) is wider than and distinct from the 6-target grid actually used in
  the frozen final run (Sec.~4.3) -- do not conflate the two.
\end{itemize}
MUST NOT SAY: that calibration numbers (Gate 1's 17-point curve) are
themselves part of the frozen 240-condition evidentiary run.
\end{draftnote}
% EVIDENCE:
% - research/EXPERIMENT_1_CALIBRATION_REPORT.md Sec.1-2

\subsection{Preregistration and Freeze Discipline}
\begin{draftnote}
PURPOSE: State explicitly that the falsification criteria were fixed before
any data existed, and what happened after the run.\\
QUESTION: What was decided before seeing results, and was anything re-tuned
after?\\
ALLOWED EVIDENCE: research/EXPERIMENT\_1\_REDESIGN\_REVIEW.md \S18, \S20.\\
KEY POINTS:
\begin{itemize}
  \item Falsification table (H1 SUPPORTED / PARTIALLY SUPPORTED / NOT
  SUPPORTED / INCONCLUSIVE) was fixed in advance, before the frozen run.
  \item The observed pattern matches the pre-registered PARTIALLY SUPPORTED
  row exactly -- not a post-hoc reinterpretation.
  \item No threshold, mechanism, retrieval-cutoff, winner-definition, or
  generator parameter was changed after the frozen run's results were known.
\end{itemize}
MUST NOT SAY: that any post-hoc analysis (D.1) altered the frozen numbers --
D.1 explains the frozen data, it does not regenerate or re-tune it.
\end{draftnote}
% EVIDENCE:
% - research/EXPERIMENT_1_REDESIGN_REVIEW.md Sec.18, Sec.20
% - research/EXPERIMENT_1_EVIDENCE_CHECKPOINT.md Sec.11 (no post-hoc tuning)

\subsection{Reproducibility}
\begin{draftnote}
PURPOSE: State exactly what a stranger can and cannot re-run, per
research/MANUSCRIPT\_ARCHITECTURE.md \S12's four-part structure.\\
QUESTION: What is public, what is synthetic, what is confidential, and where
do the artifacts live?\\
ALLOWED EVIDENCE: research/PHASE\_E\_PLAN.md Task 6; research/PUBLIC\_
RELEASE\_BOUNDARY.md.\\
KEY POINTS:
\begin{itemize}
  \item PUBLIC: the entire generator, mechanisms, perturbation model,
  calibration and analysis code (\texttt{scripts/experiments/exp1/*}) --
  offline, no API keys, no cost. Test suite (5 files, 37 tests) passing as of
  the E0 checkpoint audit.
  \item SYNTHETIC: the 240-condition dataset itself is generated, not
  collected -- re-running the committed seed manifest reproduces
  \texttt{final\_condition\_results.csv} exactly.
  \item CONFIDENTIAL: the production case study (Sec.~1.1, Sec.~6.1) -- cited
  aggregate statistics only, explicitly non-reproducible from this repository.
  \item Command: \texttt{python scripts/experiments/exp1/run\_final.py}.
\end{itemize}
MUST NOT SAY: that the production case study is reproducible in any sense.
\end{draftnote}
% EVIDENCE:
% - research/PHASE_E_PLAN.md Task 6
% - research/PUBLIC_RELEASE_BOUNDARY.md Sec.3
% - research/EXPERIMENT_1_FINAL_RESULTS.md Sec.10

% FIGURE F1 -- Experimental design and pre-registration flow.
% Purpose: show the factorial structure (6 targets x 2 lexical conditions x
%   20 seeds) and the pre-registration -> falsification-table pipeline before
%   any result is shown.
% Source artifact: research/EXPERIMENT_1_REDESIGN_REVIEW.md Sec.6, Sec.18.
% Expected variables: no axes (flow diagram); two lexical conditions as
%   parallel tracks.
% Required or optional: REQUIRED (research/MANUSCRIPT_ARCHITECTURE.md Sec.7).

% TABLE T3 -- Experimental configuration.
% Purpose: one place with every frozen parameter, for reproducibility audit.
% Source artifact: research/EXPERIMENT_1_FINAL_RESULTS.md Sec.3;
%   research/EXPERIMENT_1_EVIDENCE_CHECKPOINT.md Sec.5.
% Expected variables/rows: targets, lexical conditions, seeds, retrieval
%   cutoff, R3 thresholds, winner rule (delta, resamples, alpha), primary
%   metric, LLM status, total conditions.
% Required or optional: REQUIRED.


% =============================================================================
\section{Results}
% =============================================================================
% MOST IMPORTANT SECTION. Governing rule: (A) ADS -> mechanism accuracy and
% (B) ADS -> mechanism ranking are NEVER merged (research/CONTRIBUTION_LOCK.md
% Sec.4; research/PAPER_CONTRACT.md Sec.2 rows 3-4). Leads with the band
% structure, not the flat aggregate.
%
% NOTE ON SUBSECTION COUNT: research/MANUSCRIPT_ARCHITECTURE.md Sec.4 drafted
% Results as 6 subsections (5.1 overview/regime structure, 5.2 A, 5.3 CLEAN vs
% VARIED, 5.4 B, 5.5 mechanism-winner/R3-by-region, 5.6 uncertainty/stats).
% This file instead follows the 7-subsection breakdown given explicitly in the
% E2 kickoff prompt (5.1 completeness, 5.2 accuracy, 5.3 ranking, 5.4 R3
% threshold agreement, 5.5 ADS x representation-stability interaction, 5.6
% statistical interpretation, 5.7 summary) -- a deliberate refinement of E1's
% plan by the E2 brief's own later, more specific instruction, not an
% undocumented drift. The accuracy/ranking separation (5.2 vs 5.3) and the
% band-before-aggregate ordering (5.4) are preserved exactly either way.

\subsection{Experimental Completeness and Frozen Design}
\begin{draftnote}
PURPOSE: Establish, before any substantive finding, that the run is complete,
clean, and matches its frozen design exactly -- a reader auditing the paper
needs this before trusting anything in 5.2--5.7.\\
QUESTION: Did the experiment run as designed, completely and without error?\\
ALLOWED EVIDENCE: research/EXPERIMENT\_1\_EVIDENCE\_CHECKPOINT.md \S3-5;
research/EXPERIMENT\_1\_FINAL\_RESULTS.md \S2.\\
KEY POINTS:
\begin{itemize}
  \item 240/240 conditions succeeded, 0 failed.
  \item All 240 (seed, target, lexical) triples unique, no duplicates.
  \item Every frozen-config assertion passed (\texttt{run\_final.py} lines
  50-54: $\delta=0.02$, cutoff=75, \texttt{P\_TRANSFORM}=0.3, 20 seeds, R3
  thresholds 0.90/0.70).
  \item Full harness test suite (37 tests, 8 files -- now 5 test files per the
  current committed tree) passing.
\end{itemize}
MUST NOT SAY: nothing to forbid here -- pure integrity reporting.
\end{draftnote}
% EVIDENCE:
% - research/EXPERIMENT_1_EVIDENCE_CHECKPOINT.md Sec.3-5, Sec.11
% - research/EXPERIMENT_1_FINAL_RESULTS.md Sec.2
% - data/outputs/experiments/exp1/final/final_condition_results.csv (240 rows)

\subsection{ADS Predicts Individual Mechanism Accuracy}
\begin{draftnote}
PURPOSE: Report the positive half (6a) precisely, per mechanism, never merged
into a single "ADS predicts performance" sentence.\\
QUESTION: Does realized ADS correlate with each mechanism's own accuracy?\\
ALLOWED EVIDENCE: research/CONTRIBUTION\_LOCK.md \S4; data/outputs/
experiments/exp1/posthoc/posthoc\_analysis\_report.json.\\
KEY POINTS:
\begin{itemize}
  \item Pearson $r$(realized ADS, rules accuracy) $\approx 0.909$ (VARIED) /
  $0.959$ (CLEAN).
  \item Pearson $r$(realized ADS, retrieval accuracy) $\approx 0.948$ (VARIED)
  / $0.955$ (CLEAN).
  \item True, strongly supported, in \emph{both} lexical conditions -- report
  both conditions side by side, never only the stronger one.
  \item Explicitly correlational -- no causal claim here.
\end{itemize}
MUST NOT SAY: "ADS predicts mechanism suitability" or any sentence that
elides which mechanism this correlation is about.
\end{draftnote}
% EVIDENCE:
% - research/CONTRIBUTION_LOCK.md Sec.4 (6a)
% - data/outputs/experiments/exp1/final/final_condition_results.csv
% - research/AUDIT_REPORT.md (independent re-derivation of these r values)

% FIGURE F2 -- ADS vs. mechanism accuracy.
% Purpose: visualize 6a's correlations directly.
% Source artifact: data/outputs/experiments/exp1/final/final_condition_results.csv
% Expected variables: x = realized ADS (0.44-0.93); y = whole-set accuracy;
%   series = {rules, retrieval}; facet/color = lexical condition (CLEAN/VARIED).
% Required or optional: REQUIRED.

\subsection{ADS Does Not Predict Mechanism Ranking}
\begin{draftnote}
PURPOSE: Report the negative half (6b) as its own subsection, using the
winner-constancy result as its evidentiary basis -- never adjacent-paragraph
merged with 5.2.\\
QUESTION: Does realized ADS predict \emph{which} mechanism wins?\\
ALLOWED EVIDENCE: research/CONTRIBUTION\_LOCK.md \S4, \S6b;
research/EXPERIMENT\_1\_POSTHOC\_ANALYSIS.md \S6.\\
KEY POINTS:
\begin{itemize}
  \item Under VARIED, \texttt{empirical\_winner}=="retrieval" in 120/120
  conditions -- exceptionless, unconditional on realized ADS across its full
  observed range (0.44--0.93).
  \item Under CLEAN, \texttt{empirical\_winner}=="tie" in 120/120 conditions
  -- also exceptionless.
  \item The lexical condition is a \textbf{manipulated, controlled}
  experimental factor, not an observational one -- this is what licenses
  "governed by," not merely "correlated with," in the Discussion (Sec.~6).
  \item State plainly: because ADS is invariant to the lexical condition by
  construction (Sec.~4.4), and the winner is a constant function of the
  lexical condition, ADS structurally cannot track what determines the
  winner.
\end{itemize}
MUST NOT SAY: "higher ADS means rules is better" (reversed in this data --
PAPER\_CONTRACT.md \S3 row 9); any causal claim beyond what "manipulated,
controlled factor" licenses.
\end{draftnote}
% EVIDENCE:
% - research/CONTRIBUTION_LOCK.md Sec.4 (6b), Sec.2 step 8
% - research/EXPERIMENT_1_POSTHOC_ANALYSIS.md Sec.6
% - data/outputs/experiments/exp1/final/final_condition_results.csv

% FIGURE F4 -- Mechanism-ranking constancy across the lexical condition.
% Purpose: visualize the exceptionless 120/120 (VARIED, retrieval) vs.
%   120/120 (CLEAN, tie) winner-constancy claim.
% Source artifact: data/outputs/experiments/exp1/final/final_condition_results.csv
% Expected variables: x = realized ADS (0.44-0.93); y = rules-retrieval
%   accuracy difference; series = {CLEAN, VARIED}; zero-line marked.
% Required or optional: REQUIRED -- second most important figure in the paper.

\subsection{R3 Threshold Agreement by Realized-ADS Region}
\begin{draftnote}
PURPOSE: Report the paper's headline finding, leading with the band
structure, not the flat aggregate (per the E2 brief's explicit instruction).\\
QUESTION: How often does the frozen R3 rule agree with the empirical winner,
and does that rate depend on which realized-ADS region a condition falls in?\\
ALLOWED EVIDENCE: research/CONTRIBUTION\_LOCK.md \S2 step 7;
research/EXPERIMENT\_1\_POSTHOC\_ANALYSIS.md \S4-5.\\
KEY POINTS:
\begin{itemize}
  \item PRIMARY (per-row realized-ADS banding, matches R3's own actual
  mechanism -- CONTRIBUTION\_LOCK.md's locked wording): \textbf{100\% (32/32)}
  agreement in the realized 0.70--0.90 ADS band; \textbf{0\% (0/18)} in the
  realized $\geq$0.90 band.
  \item SECONDARY, clearly labeled as a different (also valid) slicing --
  the by-nominal-target breakdown (EXPERIMENT\_1\_FINAL\_RESULTS.md \S5): 30/30
  = 100\% at targets 0.50/0.75; 2/20 = 10\% at target 1.00. State explicitly
  \emph{why} these numbers differ slightly from the per-row band numbers
  (Sec.~4.4: nominal target and per-row realized ADS do not always coincide)
  -- do not silently switch between the two without saying which is reported.
  \item Only \emph{after} the band table: report the flat aggregate --- 32/50
  = 64.0\%, Wilson 95\% CI [50.14\%, 75.86\%], exact binomial $p=0.0649$ (not
  significant at $\alpha=0.05$) -- explicitly labeled as a statistic that
  obscures the band structure just shown, a near-cancellation of two large,
  opposite, individually highly significant effects (band $p$-values
  $1.9\times10^{-9}$ and $4.0\times10^{-4}$ respectively), not a single
  moderate effect.
  \item CLEAN contributes zero defined comparisons at any band -- all
  measurable H1 evidence comes from VARIED.
\end{itemize}
MUST NOT SAY: the 64.0\% aggregate as the paper's headline number, presented
before or without the band table; any implication that H1 was fully
confirmed (verdict is PARTIALLY\_SUPPORTED -- Sec.~7.10).
\end{draftnote}
% EVIDENCE:
% - research/CONTRIBUTION_LOCK.md Sec.2 step 7, Sec.4
% - research/EXPERIMENT_1_POSTHOC_ANALYSIS.md Sec.4-5
% - research/EXPERIMENT_1_FINAL_RESULTS.md Sec.5-6
% - research/EXPERIMENT_1_EVIDENCE_CHECKPOINT.md Sec.6-9

% FIGURE F3 -- R3 agreement by realized-ADS region.
% Purpose: the paper's headline finding -- 100% vs 0% -- shown as a single,
%   instantly legible bar chart.
% Source artifact: research/EXPERIMENT_1_POSTHOC_ANALYSIS.md Sec.4-5.
% Expected variables: x = realized-ADS band (<0.70, 0.70-0.90, >=0.90);
%   y = agreement rate (%); grouping = CLEAN vs VARIED panels/bars.
% Required or optional: REQUIRED -- the single most important figure.

% TABLE T4 -- Main results table.
% Purpose: the paper's central quantitative claim in one canonical table
%   other sections point back to.
% Source artifact: research/EXPERIMENT_1_EVIDENCE_CHECKPOINT.md Sec.6-7.
% Expected variables/rows: overall agreement (32/50, 64.0%, CI, p); band-level
%   split (32/32 vs 0/18, or 30/30 vs 2/20 by-target -- pick ONE primary
%   framing per Sec.5.4's KEY POINTS and label it); Pearson r ranges (6a).
% Required or optional: REQUIRED.

\subsection{ADS $\times$ Representation-Stability Interaction}
\begin{draftnote}
PURPOSE: Show the mechanistic pattern underlying 5.4's reversal -- the
widening rules-retrieval gap under noise as ADS rises -- without causal
overclaiming.\\
QUESTION: How does the rules-retrieval accuracy gap behave as realized ADS
increases, under each lexical condition?\\
ALLOWED EVIDENCE: research/EXPERIMENT\_1\_POSTHOC\_ANALYSIS.md \S5.\\
KEY POINTS:
\begin{itemize}
  \item Six-row band table (realized-ADS band $\times$ lexical condition
  $\rightarrow$ rules acc, retrieval acc, gap): under VARIED, the rules$-$
  retrieval gap \emph{widens} from $-0.137$ (band $<0.70$) to $-0.168$ (band
  $0.70$--$0.90$) to $-0.185$ (band $\geq0.90$) -- the opposite of what R3's
  threshold design assumes.
  \item Under CLEAN, the gap stays small and positive ($+0.005$ to $+0.006$)
  across all three bands -- near-equivalence throughout, absent noise.
  \item This table is the direct mechanical explanation for why R3's
  agreement rate is 100\% in the middle band and 0\% at the top: R3
  "succeeds" only where its ADS-driven output happens to equal the constant
  true winner.
  \item State the causal account (ADS's blindness to surface form) here only
  as a pointer to Sec.~6.4 -- do not develop the full causal argument in
  Results; Results reports the pattern, Discussion explains it.
\end{itemize}
MUST NOT SAY: a causal claim beyond "this table shows the pattern that
Discussion Sec.~6.4 explains" -- the causal account itself is INFERRED,
post-hoc (Sec.~7 Limitations), and belongs in Discussion, not Results.
\end{draftnote}
% EVIDENCE:
% - research/EXPERIMENT_1_POSTHOC_ANALYSIS.md Sec.5
% - research/CONTRIBUTION_LOCK.md Sec.2 step 8 (pointer only, not the full account)

% TABLE T5 -- Mechanism-winner behavior by ADS region.
% Purpose: exact numbers for the widening-gap pattern -- a figure would lose
%   precision a reviewer checking this claim would want.
% Source artifact: research/EXPERIMENT_1_POSTHOC_ANALYSIS.md Sec.5 (the
%   six-row table, reproduced verbatim, not re-derived).
% Expected variables/rows: realized-ADS band x lexical condition -> n, rules
%   acc (mean), retrieval acc (mean), rules-retrieval gap, rules/retrieval/tie
%   win counts.
% Required or optional: REQUIRED.

\subsection{Statistical Interpretation}
\begin{draftnote}
PURPOSE: Place the CI/binomial-test machinery last, after the substantive
finding it supports, not before.\\
QUESTION: How confident should a reader be in the numbers reported above?\\
ALLOWED EVIDENCE: research/EXPERIMENT\_1\_EVIDENCE\_CHECKPOINT.md \S6-9.\\
KEY POINTS:
\begin{itemize}
  \item Both individual bands are far more statistically significant than the
  aggregate ($p=1.9\times10^{-9}$ and $p=4.0\times10^{-4}$ vs.\ the
  aggregate's $p=0.0649$).
  \item $n_{\mathrm{defined}}$ is small in the two informative bands (30 and
  20 conditions) -- enough for the binomial tests to be decisive, not enough
  to subdivide further.
  \item Restate that $\delta=0.02$ (frozen at Gate 3) is the specific reason
  CLEAN contributes zero defined comparisons -- a property of the frozen
  definition, not re-examined here.
\end{itemize}
MUST NOT SAY: any relaxation of $\alpha=0.05$ or re-interpretation of the
frozen data to reach a different verdict.
\end{draftnote}
% EVIDENCE:
% - research/EXPERIMENT_1_EVIDENCE_CHECKPOINT.md Sec.6-9
% - research/EXPERIMENT_1_REDESIGN_REVIEW.md Sec.18 (falsification criteria)

\subsection{Summary of Findings}
\begin{draftnote}
PURPOSE: One short paragraph restating 5.2--5.6 without adding anything new --
a bridge into Discussion, not a new claim.\\
QUESTION: What, in total, did Results establish?\\
ALLOWED EVIDENCE: This section's own subsections above only.\\
KEY POINTS:
\begin{itemize}
  \item 6a holds (5.2); 6b holds (5.3); the mechanism is the ADS$\times$
  representation-stability interaction (5.5); the aggregate alone
  under-communicates the structure (5.4).
\end{itemize}
MUST NOT SAY: any claim not already stated in 5.1--5.6 above.
\end{draftnote}


% =============================================================================
\section{Discussion}
% =============================================================================
% The E2 brief's eight questions, in order, plus one additional subsection
% ("The Production Case Study in Light of These Results", inserted after the
% first question) required by research/MANUSCRIPT_ARCHITECTURE.md Sec.6's
% explicit instruction for a "clearly labeled short subsection" positioning
% the production observation relative to 6a/6b -- flagged as a required fix
% by research/MANUSCRIPT_SKELETON_AUDIT.md, added here rather than left to E3.

\subsection{What the Experiment Supports}
\begin{draftnote}
PURPOSE: State the 6a/6b synthesis, exactly as locked.\\
QUESTION: What did we learn?\\
ALLOWED EVIDENCE: research/CONTRIBUTION\_LOCK.md \S6 (Synthesis sentence).\\
KEY POINTS:
\begin{itemize}
  \item Historical decision consistency is informative about classification-
  mechanism \emph{difficulty}, not mechanism \emph{ranking}, when -- as
  observed here -- ranking is governed by a representation-stability property
  the consistency signal does not observe.
  \item This is a narrower, more specific, empirically falsifiable-and-
  partly-falsified refinement of the original hypothesis, not a confirmation
  of it.
\end{itemize}
MUST NOT SAY: anything stronger than CONTRIBUTION\_LOCK.md \S6's synthesis
sentence.
\end{draftnote}
% EVIDENCE:
% - research/CONTRIBUTION_LOCK.md Sec.6

\subsection{The Production Case Study in Light of These Results}
\begin{draftnote}
PURPOSE: The clearly-labeled short subsection research/MANUSCRIPT\_ARCHITECTURE.md
\S6 requires -- position the production observation relative to this
experiment's actual findings, once, explicitly, rather than leaving the
relationship implicit.\\
QUESTION: How does the production "R3 flip" observation (Sec.~1.1) relate to
what this experiment actually found?\\
ALLOWED EVIDENCE: research/EVIDENCE\_BASELINE.md \S1; research/r3\_threshold\_
analysis.md.\\
KEY POINTS:
\begin{itemize}
  \item The production observation is \emph{consistent with} 6a (ADS did
  correlate with something real in production) but was \textbf{never itself a
  controlled test of 6b} -- production never ran a lexical-noise sweep.
  \item State plainly: "this experiment does not confirm the production
  observation, it investigates the more general question the observation
  raised."
  \item Production ADS figures remain pre-A5-fix, "likely understated,
  unverified" (research/EVIDENCE\_BASELINE.md \S1) -- repeat the caveat here,
  do not assume Sec.~1.1 already carries it far enough for a reader who skips
  ahead to Discussion.
\end{itemize}
MUST NOT SAY: that this experiment confirms, validates, or independently
tests the production observation; that the production case study underwent
the same lexical-noise manipulation as Experiment 1 (research/MANUSCRIPT\_
ARCHITECTURE.md \S6, explicit prohibition).
\end{draftnote}
% EVIDENCE:
% - research/EVIDENCE_BASELINE.md Sec.1
% - research/r3_threshold_analysis.md
% - research/MANUSCRIPT_ARCHITECTURE.md Sec.6 (production case study placement)

\subsection{What the Original Hypothesis Got Right}
\begin{draftnote}
PURPOSE: Credit H1's weaker, revised form honestly -- it is not cleanly
confirmed or rejected.\\
QUESTION: What did the original hypothesis get right?\\
ALLOWED EVIDENCE: research/EXPERIMENT\_1\_REDESIGN\_REVIEW.md \S18;
research/EXPERIMENT\_1\_EVIDENCE\_CHECKPOINT.md \S9.\\
KEY POINTS:
\begin{itemize}
  \item H1's weaker form ("better-than-chance agreement," not "monotonic
  dominance in every band") is exactly why a PARTIALLY\_SUPPORTED verdict was
  even a coherent, pre-registered possible outcome.
  \item The point estimate (64.0\%) is directionally above chance in the only
  slice where agreement is measurable (VARIED), and one full band shows
  perfect, highly significant agreement.
\end{itemize}
MUST NOT SAY: that this amounts to "H1 confirmed."
\end{draftnote}
% EVIDENCE:
% - research/EXPERIMENT_1_REDESIGN_REVIEW.md Sec.18
% - research/EXPERIMENT_1_EVIDENCE_CHECKPOINT.md Sec.9

\subsection{What the Original Hypothesis Got Wrong}
\begin{draftnote}
PURPOSE: State the falsified unconditional form plainly.\\
QUESTION: What did the original hypothesis get wrong?\\
ALLOWED EVIDENCE: research/CONTRIBUTION\_LOCK.md \S3 (C2b downgrade).\\
KEY POINTS:
\begin{itemize}
  \item The \emph{unconditional} form ("ADS thresholds alone select the right
  mechanism") is falsified, exceptionlessly, in the realized $\geq$0.90 band
  (0/18).
  \item This is exactly the band R3 is most confident in (highest
  consistency $\rightarrow$ strongest recommendation for rules) -- the
  failure is not at the margins, it is at R3's own point of maximum
  confidence.
\end{itemize}
MUST NOT SAY: that this falsification is minor or an edge case -- it is
exceptionless and occurs at R3's highest-confidence point.
\end{draftnote}
% EVIDENCE:
% - research/CONTRIBUTION_LOCK.md Sec.3 (C2b), Sec.4
% - research/EXPERIMENT_1_FINAL_RESULTS.md Sec.6, Sec.8

\subsection{Why Consistency Predicts Difficulty but Not Ranking}
\begin{draftnote}
PURPOSE: Give the mechanistic account, explicitly labeled as INFERRED and
post-hoc.\\
QUESTION: Why is ADS informative about difficulty but insufficient for
ranking?\\
ALLOWED EVIDENCE: research/CONTRIBUTION\_LOCK.md \S2 step 8.\\
KEY POINTS:
\begin{itemize}
  \item ADS is computed on the stable \texttt{product\_code}, structurally
  blind to the perturbable surface string.
  \item Ranking, in this experiment, is a constant function of the lexical
  condition -- a variable ADS cannot see by construction.
  \item State explicitly: this account is INFERRED from exhaustive-but-post-
  hoc code inspection and matching, not a second, independently-designed
  confirmatory experiment.
\end{itemize}
MUST NOT SAY: "we prove" or "we demonstrate causally" -- use "the evidence is
consistent with," "the exhaustive pattern supports," or equivalent hedged
language throughout this subsection.
\end{draftnote}
% EVIDENCE:
% - research/CONTRIBUTION_LOCK.md Sec.2 step 8
% - research/EXPERIMENT_1_DATA_DICTIONARY.md ("What realized ADS is, precisely")

\subsection{Representation Stability as an Uncaptured Factor}
\begin{draftnote}
PURPOSE: Discuss what the widening-gap pattern (5.5) reveals about
representation stability, scoped tightly to this perturbation model.\\
QUESTION: What does representation stability reveal here?\\
ALLOWED EVIDENCE: research/EXPERIMENT\_1\_POSTHOC\_ANALYSIS.md \S5-6.\\
KEY POINTS:
\begin{itemize}
  \item Retrieval's accuracy cost stays comparatively flat under noise while
  rules' cost grows with ADS -- the perturbation degrades an exact-match
  mechanism's precondition (identical surface strings) in a way a fuzzy
  mechanism is partly robust to.
  \item Framed as a property of \emph{this} perturbation model and \emph{this}
  retrieval implementation (rapidfuzz) -- not fuzzy/embedding retrieval
  generally (PAPER\_CONTRACT.md \S6).
\end{itemize}
MUST NOT SAY: that this generalizes to other retrieval implementations
(e.g.\ embedding-based) or other noise models -- untested, named only as
future work (Sec.~8).
\end{draftnote}
% EVIDENCE:
% - research/EXPERIMENT_1_POSTHOC_ANALYSIS.md Sec.5-6
% - research/PAPER_CONTRACT.md Sec.6

\subsection{Relationship to Algorithm Selection and Meta-Learning}
\begin{draftnote}
PURPOSE: Situate the finding against the Related Work lineage, as a
refinement, not a contradiction.\\
QUESTION: How does this relate to Algorithm Selection / meta-learning?\\
ALLOWED EVIDENCE: research/CONTRIBUTION\_LOCK.md \S5 row 2;
research/literature/terminology\_map.md.\\
KEY POINTS:
\begin{itemize}
  \item This paper refines, not contradicts, the algorithm-selection/meta-
  learning lineage \citep{rice1976,smithmiles2009,barbudo2023} by identifying
  a specific failure mode of a single-feature, label-consistency-only
  design-time selector.
\end{itemize}
MUST NOT SAY: that this literature is wrong, superseded, or that this paper's
finding is a bigger challenge to it than the narrow failure mode described.
\end{draftnote}
% EVIDENCE:
% - research/CONTRIBUTION_LOCK.md Sec.5 row 2
% - research/literature/terminology_map.md

\subsection{What Practitioners Should NOT Infer}
\begin{draftnote}
PURPOSE: A dedicated, explicit paragraph -- do not leave this implicit.\\
QUESTION: What should practitioners NOT infer from this experiment?\\
ALLOWED EVIDENCE: research/PAPER\_CONTRACT.md \S3, \S6.\\
KEY POINTS:
\begin{itemize}
  \item Do not infer that ADS-style consistency signals are useless (6a
  still holds).
  \item Do not infer this generalizes to other noise models or retrieval
  implementations.
  \item Do not infer that a two-feature fix exists (none was built -- Sec.~8).
  \item Do not infer that production data confirms this finding (it does
  not, and never ran this experiment -- Sec.~1.1, Sec.~7.7).
\end{itemize}
MUST NOT SAY: nothing -- this subsection IS the "must not say" list, rendered
for the reader rather than only for the drafting process.
\end{draftnote}
% EVIDENCE:
% - research/PAPER_CONTRACT.md Sec.3, Sec.6

\subsection{Implications for Future Selector Design}
\begin{draftnote}
PURPOSE: Name the natural next step, as a direction only, with the explicit
guardrail against presenting it as built.\\
QUESTION: What future research naturally follows?\\
ALLOWED EVIDENCE: research/CONTRIBUTION\_LOCK.md \S10.\\
KEY POINTS:
\begin{itemize}
  \item A decision rule conditioning on both ADS and a measured (not assumed)
  representation-stability signal -- the direct next hypothesis this
  experiment motivates.
  \item State in future tense only, with no pseudocode, no proposed threshold
  values, no partial results.
\end{itemize}
MUST NOT SAY: any implementation detail beyond "condition on both signals" --
doing so would cross from "naming a direction" into "presenting an untested
method as if scoped" (PAPER\_CONTRACT.md \S3 row 12).
\end{draftnote}
% EVIDENCE:
% - research/CONTRIBUTION_LOCK.md Sec.10
% - research/PAPER_CONTRACT.md Sec.3 row 12


% =============================================================================
\section{Limitations}
% =============================================================================
% Every item from research/CONTRIBUTION_LOCK.md Sec.9 and
% research/PAPER_CONTRACT.md Sec.6, undiluted. Negative finding not buried:
% Sec.7.10 leads with the honest PARTIALLY_SUPPORTED verdict.

\subsection{Synthetic Generator Scope}
\begin{draftnote}
KEY POINTS:
\begin{itemize}
\item Single synthetic generator family; 240 conditions replicate a seeded
RNG, not an external population.
\end{itemize}
\end{draftnote}
% EVIDENCE: research/CONTRIBUTION_LOCK.md Sec.9

\subsection{Lexical Perturbation Model}
\begin{draftnote}
KEY POINTS:
\begin{itemize}
\item \texttt{p\_transform=0.3}, five fixed transform types, pilot-tuned to a
target corruption-share band -- not derived from measured real-world
OCR/typo error rates.
\end{itemize}
\end{draftnote}
% EVIDENCE: research/CONTRIBUTION_LOCK.md Sec.9; research/EXPERIMENT_1_REDESIGN_REVIEW.md Sec.22

\subsection{Single Application Domain}
\begin{draftnote}
KEY POINTS:
\begin{itemize}
\item The motivating case study is one domain (Romanian fiscal-document /
invoice GL-account classification); the experiment itself is domain-agnostic
by construction, but has not been run against a second domain's data.
\end{itemize}
\end{draftnote}
% EVIDENCE: research/CONTRIBUTION_LOCK.md Sec.8

\subsection{Only Two Mechanisms Compared}
\begin{draftnote}
KEY POINTS:
\begin{itemize}
\item Exact-match rules vs.\ rapidfuzz-based retrieval only -- not embeddings,
not the shipped multi-tier production cascade.
\end{itemize}
\end{draftnote}
% EVIDENCE: research/CONTRIBUTION_LOCK.md Sec.8

\subsection{LLM Exclusion}
\begin{draftnote}
KEY POINTS:
\begin{itemize}
\item The LLM mechanism was excluded from H1 for a documented, principled
reason (the synthetic product string carries no signal predictive of the true
label) -- not evaluated at all in Experiment 1.
\end{itemize}
\end{draftnote}
% EVIDENCE: research/EXPERIMENT_1_REDESIGN_REVIEW.md Sec.10; research/CONTRIBUTION_LOCK.md Sec.9

\subsection{Production Confidentiality}
\begin{draftnote}
KEY POINTS:
\begin{itemize}
\item Production ADS figures (91.2\%/0.847/0.964) were generated by the
pre-A5-fix script version and remain "likely understated, unverified" -- not
re-run against corrected code, since no production data exists in this
repository.
\end{itemize}
\end{draftnote}
% EVIDENCE: research/EVIDENCE_BASELINE.md Sec.1; research/CONTRIBUTION_LOCK.md Sec.9

\subsection{Absence of Independent Production Validation}
\begin{draftnote}
KEY POINTS:
\begin{itemize}
\item Production never ran a lexical-noise sweep at all; only two single-run
data points feed the R3-flip narrative -- production data does not
independently validate this experiment's finding, and this must never be
implied.
\end{itemize}
\end{draftnote}
% EVIDENCE: research/CONTRIBUTION_LOCK.md Sec.7 (rejected claim), Sec.8

\subsection{No Two-Feature Selector}
\begin{draftnote}
KEY POINTS:
\begin{itemize}
\item No decision rule conditioning on both ADS and a representation-
stability signal was designed, prototyped, or tested anywhere in this
repository -- named only as future work (Sec.~8.1).
\end{itemize}
\end{draftnote}
% EVIDENCE: research/CONTRIBUTION_LOCK.md Sec.5 row 3-4, Sec.10

\subsection{Limited Generalization}
\begin{draftnote}
KEY POINTS:
\begin{itemize}
\item The realized-ADS ceiling ($\sim$0.91, from the fixed cross-company-
alignment nuisance parameter) means the "deep rules-first" region
($\geq$0.93) was never reachable or tested -- untested, not merely
unfavorable.
\item No deployment or generalization claim: the finding describes this
experiment, not a recommendation for how to build classification systems in
general.
\end{itemize}
\end{draftnote}
% EVIDENCE: research/CONTRIBUTION_LOCK.md Sec.8-9; research/PAPER_CONTRACT.md Sec.6

\subsection{H1 Only Partially Supported}
\begin{draftnote}
PURPOSE: Lead this group with the honest verdict -- do not bury it after a
list of scope caveats that could read as excuse-making.\\
KEY POINTS:
\begin{itemize}
\item H1 (revised) is only \emph{partially} supported, matching the pre-
registered PARTIALLY SUPPORTED row exactly -- stated as the honest verdict,
not softened toward either "confirmed" or "refuted."
\item D.1's causal account (ADS's blindness to surface form) is a well-
evidenced, exhaustive, but \emph{post-hoc} explanation of already-frozen
data, not a fresh prospective confirmation.
\item Retrieval coverage is 1.0 in all 240 conditions (never abstains at
cutoff=75); rules coverage is always $<1.0$ -- a real asymmetry, part of the
explanation, not normalized away.
\end{itemize}
MUST NOT SAY: any hedge that softens this verdict into a strength.
\end{draftnote}
% EVIDENCE:
% - research/EXPERIMENT_1_REDESIGN_REVIEW.md Sec.18 (PARTIALLY SUPPORTED row)
% - research/EXPERIMENT_1_EVIDENCE_CHECKPOINT.md Sec.9
% - research/CONTRIBUTION_LOCK.md Sec.9

% TABLE T6 (optional) -- Limitations/reproducibility summary.
% Purpose: one row per limitation category above, tagged reproducible?
% Source artifact: this section's own subsections.
% Expected variables/rows: 10 rows, one per subsection 7.1-7.10.
% Required or optional: OPTIONAL -- include only if the prose list above
%   reads as too long to scan at E3; defer the decision, do not pre-commit.


% =============================================================================
\section{Future Work}
% =============================================================================
% Every item below is named, not built. All explicitly hedged as untested.

\subsection{Representation-Stability-Aware Selectors}
\begin{draftnote}
KEY POINTS:
\begin{itemize}
\item A decision rule that conditions mechanism selection on both realized
ADS and a measured (not assumed) representation-stability signal.
\end{itemize}
MUST NOT SAY: anything implying this was designed, prototyped, or partially
tested -- named only.
\end{draftnote}
% EVIDENCE: research/CONTRIBUTION_LOCK.md Sec.10 (item 1)

\subsection{Testing Additional Domains}
\begin{draftnote}
KEY POINTS:
\begin{itemize}
\item Re-running the corrected pipeline against a second domain's data, to
test whether the method's design-agnosticism (Sec.~3.1) actually transfers.
\end{itemize}
\end{draftnote}
% EVIDENCE: TECHNICAL_REPORT.md Sec.4 (domain-transfer limitation, framing only)

\subsection{Additional Mechanisms}
\begin{draftnote}
KEY POINTS:
\begin{itemize}
\item Whether the specific pattern (retrieval's noise-cost staying flat while
rules' noise-cost grows with ADS) is specific to rapidfuzz-style token
similarity or holds for other retrieval implementations (e.g.\ embedding-based
retrieval).
\end{itemize}
\end{draftnote}
% EVIDENCE: research/CONTRIBUTION_LOCK.md Sec.10 (item 2); research/EXPERIMENT_1_REDESIGN_REVIEW.md Sec.9

\subsection{Real OCR/Noise Distributions}
\begin{draftnote}
KEY POINTS:
\begin{itemize}
\item A real (not synthetic) noise/OCR-error model calibrated against an
actual document corpus, to test whether \texttt{p\_transform=0.3}'s pattern
generalizes beyond this pilot-tuned synthetic stand-in.
\end{itemize}
\end{draftnote}
% EVIDENCE: research/CONTRIBUTION_LOCK.md Sec.10 (item 3)

\subsection{Independent Datasets}
\begin{draftnote}
KEY POINTS:
\begin{itemize}
\item Re-running this experiment's design against an independently generated
dataset (not merely a different seed of the same generator) to test whether
the 6a/6b finding is a property of this generator family specifically.
\end{itemize}
\end{draftnote}
% EVIDENCE: research/CONTRIBUTION_LOCK.md Sec.9 (single-generator-family limitation)

\subsection{Broader Architecture-Selection Experiments}
\begin{draftnote}
KEY POINTS:
\begin{itemize}
\item Experiment 2 (baseline comparison against simpler pipelines, C6) and
Experiment 3 (feedback-loop before/after measurement, C5) -- recommended for
upgrading currently WEAK claims, but explicitly not required for this paper's
narrower, locked contribution.
\end{itemize}
MUST NOT SAY: that either experiment is planned, scheduled, or has begun --
neither is (research/RESEARCH\_GPS.md, "DO NOT CHASE" list).
\end{draftnote}
% EVIDENCE: research/CONTRIBUTION_LOCK.md Sec.10 (item 4); research/RESEARCH_GPS.md


% =============================================================================
\section{Conclusion}
% =============================================================================
\begin{draftnote}
PURPOSE: One paragraph. The highest-risk section in the paper for claim
inflation -- no new claim may appear here that was not already made in
Results/Discussion.\\
QUESTION: What is genuinely new/useful, in one paragraph?\\
ALLOWED EVIDENCE: research/CONTRIBUTION\_LOCK.md \S11.B/\S11.C, verbatim.\\
THE SKELETON REQUIRES THE FINAL PROSE TO CONTAIN, IN ORDER:
\begin{enumerate}
\item WHAT WAS TESTED: a pre-registered, 240-condition synthetic factorial
experiment comparing exact-match rules against fuzzy retrieval under a
controlled lexical perturbation.
\item WHAT WAS SUPPORTED: realized ADS strongly predicts each mechanism's own
accuracy (6a).
\item WHAT WAS FALSIFIED/LIMITED: ADS does not predict which mechanism wins;
the unconditional form of the original hypothesis is falsified exceptionlessly
in the realized $\geq$0.90 band (6b); H1 overall is PARTIALLY\_SUPPORTED, not
confirmed.
\item WHAT THE MAIN LESSON IS: historical consistency is informative about
mechanism difficulty, not mechanism ranking, when ranking is governed by a
representation-stability property the signal does not observe.
\item WHAT REMAINS OPEN: whether a two-feature (consistency + representation-
stability) selector would perform better -- named, not built or tested here.
\end{enumerate}
MUST NOT SAY: ``ADS solves architecture selection'' or any equivalent
closing claim. The paper's actual closing claim is a precise, evidenced
\emph{boundary} on what a historical-consistency signal can and cannot tell a
system designer -- not a solution, not a validated method, not a new metric or
architecture.
\end{draftnote}
% EVIDENCE:
% - research/CONTRIBUTION_LOCK.md Sec.11.B, Sec.11.C
% - research/PAPER_CONTRACT.md Sec.3 row 16 (Formulation #4, rejected)


% =============================================================================
\section*{Reproducibility Statement}
% =============================================================================
\begin{draftnote}
PURPOSE: A short, separate, citable statement duplicating Sec.~4.17's content
in venue-expected form (many venues expect a dedicated Reproducibility
Statement near the end of the paper, distinct from the in-body subsection).\\
KEY POINTS:
\begin{itemize}
\item Reuse Sec.~4.17's four-part structure verbatim; do not re-derive.
\item Link, do not bundle, the raw \texttt{final\_condition\_results.csv} in
the arXiv source package (research/PUBLIC\_RELEASE\_BOUNDARY.md \S3.B) --
point to the public GitHub repository instead.
\end{itemize}
\end{draftnote}
% EVIDENCE:
% - research/MANUSCRIPT_ARCHITECTURE.md Sec.12
% - research/PUBLIC_RELEASE_BOUNDARY.md Sec.3.B


% =============================================================================
% REFERENCES
% =============================================================================
% Do not perform new literature research (per the E2 brief). Every \citep{}
% key above resolves against manuscript/references.bib, itself built only
% from research/literature/citation_ledger.csv and
% research/literature/ads_metric_prior_art.csv VERIFIED rows. \nocite{*} forces
% every bib entry to render in the compiled reference list even before every
% key is cited in body prose, so an auditor can check the full candidate
% reference set against the ledger without waiting for E3 prose.
\nocite{*}
\bibliographystyle{plainnat}
\bibliography{references}

\end{document}
